\UseRawInputEncoding
\chapter{Numeri reali}


\begin{esercizio}
Motivare per quale motivo le seguenti scritture sono prive di significato: 
\[
\log_1 7 \qquad \big (\log_{1/2} 3\big )^\pi  \qquad 
\]
\end{esercizio}
\textit{Soluzione.} 
\begin{itemize}
	\item Soltanto $\log_1 1$ � una forma definita, anche se uguale a qualsiasi $k$ reale in quanto � vera:
	\[\log_1 1=k \quad \Leftrightarrow \quad 
	1^k=1 \qquad \forall \, k 
	\]
	\item Le potenze con esponente irrazionale sono definite soltanto se la base � positiva. 
\end{itemize}

\begin{esercizio}
	Stabilire per quali valori reali di $\lambda$ � definita, nel campo reale, l'espressione: 
	\[
	\mu=\log_{\lambda^2-2\lambda} \big (\log_3 \lambda)
	\]
\end{esercizio}
\textit{Soluzione}. Si osserva che, nel caso in cui la base � uguale a $1$, sono ammissibili soltanto i valori di $\lambda$ per cui anche l'argomento � uguale a $1$. Si osserva che ci� non accade: 
\[
\begin{cases} 
	\lambda^2-2\lambda=1\\
\log_3 \lambda=1
\end{cases} 
\Rightarrow \quad 
\begin{cases}
		\lambda^2-2\lambda\big |_{\lambda=3}\neq 1\\
	\lambda=3
\end{cases}
\]
Segue che la base deve essere diversa da $1$. Inoltre 
argomenti e basi dei logaritmi devono essere positivi:
\[
\begin{cases} 
\lambda>0\\
\log_3 \lambda>0\\
\lambda^2-2\lambda>0\\
\lambda^2-2\lambda\neq 1
\end{cases}
\Rightarrow \quad 
\begin{cases}
	\lambda>2\\
	\lambda\neq 1+\sqrt{2}
\end{cases}
\]

\begin{esercizio}
	Stabilire per quali valori di $x$ � reale il seguente numero $\alpha$:
	\[
	\alpha=(x-1)^{x-2}
	\]
\end{esercizio}
\textit{Soluzione}. Si procede per casi: 
\begin{itemize}
	\item Se $x>1$ allora $\alpha$ � reale. 
	\item Se $x=1$ allora si ha $\alpha=0^{-1}$, scrittura priva di significato.
	\item Se $x<1$:
	\begin{itemize}
		\item L'operazione non � definita per $x$ irrazionale.
		\item Se l'esponente � razionale allora $\alpha$ � reale se $x-2$ � un numero razionale con denominatore dispari ($\alpha$ � un radicale di indice dispari).  
	\end{itemize} 
\end{itemize}

\begin{esercizio}
	Stabilire se l'insieme dei numeri irrazionali � chiuso rispetto alle usuali operazioni di addizione e moltiplicazione. 
\end{esercizio}
\textit{Soluzione}. L'insieme degli irrazionali non � chiuso n� rispetto alla somma n� rispetto al prodotto.  \' E quindi sufficiente giustificare l'affermazione fornendo dei controesempi: 
\begin{align*}
	&\begin{cases} 
		\alpha=\sqrt{2}\in \mathcal{I}\\
		\beta=2-\sqrt{2}\in \mathcal{I}
		\end{cases} \Rightarrow \quad 
	\alpha+\beta=2\not \in \mathcal{I}
		&\begin{cases}
		\alpha=\sqrt{2}\in \mathcal{I}\\
		\beta=\sqrt{8}\in \mathcal{I}
	\end{cases}
	\Rightarrow \quad 
	\alpha\cdot \beta=4\not \in \mathcal{I}
\end{align*}

\begin{esercizio}
	Dimostrare che $\sin \frac{\pi}{18}$ � irrazionale. 
\end{esercizio}
\textit{Soluzione}. Si sfrutta una identit� goniometrica relativa a $\sin(3x)$:
\begin{align*}
	\sin 3x&=\sin (2x+x)=\sin 2x\cos x+\cos 2x\sin x\\
	&=2\sin x\, (1-\sin^2 x) +(1-2\sin^2 x)\, \sin x
	=3\sin x-4\sin^3 x
\end{align*}
Posto $x=\pi/18$:
\[
3x=\frac{\pi}{6} 
\quad \Rightarrow \quad 
\sin \frac{\pi}{6}=3\sin \frac{\pi}{18}-4\sin^3 \frac{\pi}{18}
\]
Posto inoltre $S=\sin \frac{\pi}{18}$:
\[
\frac{1}{2}=3S-4S^3 \quad \Rightarrow \quad 
8S^3-6S+1=0
\]
Per la propriet� dimostrata nell'esercizio \ref{equazrazionale}, una soluzione $S$ razionale sarebbe  del tipo: 
\[
\pm 1 \qquad \pm \frac{1}{2} \qquad \pm \frac{1}{4} \qquad \pm \frac{1}{8}
\]
Siccome l'equazione � di grado dispari e a coefficienti reali, esiste una soluzione reale. 
Siccome nessuno dei precedenti valori � soluzione, la soluzione $S$ � irrazionale. 
Si conclude che $\sin \frac{\pi}{18}$ � irrazionale. 

\begin{esercizio}
	Dimostrare che $\log_{10} 2$ � irrazionale. 
\end{esercizio}
\textit{Soluzione}. Per assurdo, se fosse razionale, esisterebbero gli interi $p,q$ primi tra loro, con $q\neq 0$:
\[
\log_{10} 2=\frac{p}{q} 
\quad \Rightarrow \quad 
2=10^{p/q} 
\quad \Rightarrow \quad 
2^q=10^p=2^p\cdot 5^p
\]
Da cui un assurdo in quanto si contraddice l'unicit� della scomposizione in fattori primi. \\
\newline
Si osserva che se $\log_{10} \alpha$ � razionale allora $\alpha$ � una potenza di $10$:
\[
\log_{10} \alpha=\frac{p}{q} \quad \Rightarrow \quad 
\alpha=10^{p/q} \quad \Rightarrow \quad 
\alpha^q=10^p
\]